Efficient and precise retrieval is a critical research problem in natural language processing, with a wide range of applications such as semantic search~\cite{semantic_search}, open domain question answering~\cite{chen-yih-2020-open,zhu2021retrieving} and Retrieval Augmented Generation(RAG)~\cite{lewis2020retrieval} to name a few.
\eat{There have been significant advancements in retriever paradigms, evolving from traditional lexical retrievers like BM25~\cite{robertson2009probabilistic}, which utilizes sparse inverted indexes, to modern dense embedding-based retrievers such as DPR~\cite{Karpukhin2020DensePR} and ColBERT~\cite{Khattab2020ColBERTEA}, which rely on dedicated VectorDBs~\cite{johnson2019billion} for indexing and retrieval. A more recent approach is Differential Search Indexing (DSI)\cite{Tay2022TransformerMA}, which introduces a sequence-to-sequence model that maps queries directly to document identifiers, thereby embedding corpus knowledge within the model parameters.} With the rise of RAG applications, precise retrieval has become increasingly important, particularly in curbing  hallucination in generative models. \eat{Consequently, a central research focus across all retriever paradigms is the development of more effective indexing strategies to enhance retrieval quality} We can broadly classify all retrieval techniques into three main categories:

\textbf{Dense Retrieval}~\cite{Cai2021SemanticMF,Karpukhin2020DensePR} represent a query and document using a dense vector and compute the similarity based on the distance between their vectors. DPR~\cite{Karpukhin2020DensePR} uses dual encoder, ColBERT ~\cite{Khattab2020ColBERTEA} represents every token in the query and document using dense vector and performs late interaction to score query and document pairs. Recently,~\citet{sarthi2024raptor} introduced RAPTOR, a method for building a hierarchical tree by recursively embedding, clustering, and summarizing content from lengthy documents to retrieve relevant text at various levels of abstraction. Our approach differs in two key aspects: while RAPTOR uses dense retrievers for indexing and retrieval, we adopt model-based indexing. Additionally, we enhance model-based indexing by leveraging table-of-contents ({\toc}) structures, unlike RAPTOR, which relies solely on dense retrievers to access information at multiple levels of abstraction.

\textbf{Learned Sparse Retrieval} systems use sparse vector representation for a query and document. Unlike traditional sparse retrieval methods like BM25\cite{robertson2009probabilistic}, which rely on exact token matches, these models learn to represent queries and documents in a sparse high-dimensional space, allowing for efficient lexical matching. SPLADE~\cite{Hai2023CoSPLADECS} (Sparse Lexical and Expansion Model for Information Retrieval) is a prominent example of such systems. Several other models have been developed to enhance sparse retrieval, including SPLADEv2 \cite{formal2021splade}, DeepCT\cite{dai2019context}, uniCOIL\cite{lin2021few} and DeepImpact~\cite{basnet2024deeperimpact}.

\textbf{Model-based Indexing}~\cite{Metzler2021RethinkingSM} such as Differential Search Index (DSI)~\cite{Tay2022TransformerMA} embed knowledge of the entire corpus directly into its model parameters -- greatly simplifying the retrieval process. {\sys} takes this paradigm one step further by learning to exploit a rich hierarchal and semantic global structure within the corpus.